%! Logarithmic Function Context and Data Modeling — Unit 5, Lesson 5.7 — Lesson Plan
\documentclass[10pt]{article}
\usepackage{precalculus-article}
\usepackage{precalculus-boxes}
\usepackage{graphicx}
\graphicspath{{images/}}
\newcommand{\TallMath}[1]{$\displaystyle #1\rule[-1.4em]{0pt}{3.2em}$}
\newcommand{\CourseName}{Precalculus}
\newcommand{\SchoolYear}{2026--2027}
\newcommand{\MeetingLength}{55 minutes}
\newcommand{\UnitNumberName}{Unit 5: Logarithmic Functions \quad}
\newcommand{\LessonNumberName}{Lesson 5.7: Logarithmic Function Context and Data Modeling}

\begin{document}
\begin{center}
    {\Large\bfseries \CourseName: \SchoolYear \\
    {\small\bfseries \UnitNumberName \LessonNumberName}}
\end{center}

\begin{tcolorbox}[colback=sky, colframe=navy, boxrule=0.9pt, arc=2mm,
    left=2mm, right=2mm, top=1.5mm, bottom=1.5mm]
    \textbf{Primary Objective:} Students construct logarithmic function models from contextual
    data and two input-output pairs, apply transformations of $f(x) = a\log_b x$ to fit
    real-world situations, use natural logarithm models for real phenomena, and use models to
    predict values for the dependent variable.
    \hfill\textit{\small Skills: AP 1.C, AP 3.B \quad LO 2.14.A}
\end{tcolorbox}

\renewcommand{\arraystretch}{1.6}
\begin{skillbox}[Priority Ideas \& Skills]{goldbox}
\begin{tabularx}{\linewidth}{@{}X X@{}}
\textbf{AP Skills} &
\textbf{Key Understandings} \\
\midrule
\textbf{AP 1.C} --- Construct new functions using transformations and modeling. &
\begin{itemize}[leftmargin=1.2em,itemsep=2pt,topsep=0pt]
  \item Logarithmic models arise when inputs change proportionally over equal-length output
        intervals. (EK 2.14.A.1)
  \item A model can be built from a proportion + real zero, or from two input-output pairs
        via a system of equations. (EK 2.14.A.2)
  \item Transformations of $f(x)=a\log_b(x-h)+k$ encode context characteristics.
        (EK 2.14.A.3)
  \item The natural logarithm is especially useful for modeling real phenomena.
        (EK 2.14.A.5)
\end{itemize} \\
\textbf{AP 3.B} --- Apply numerical results in context. &
\begin{itemize}[leftmargin=1.2em,itemsep=2pt,topsep=0pt]
  \item Use the model to predict values of the dependent variable and interpret results
        in context. (EK 2.14.A.6)
  \item Technology (logarithmic regression) can construct models for data sets; results
        must be interpreted contextually. (EK 2.14.A.4)
\end{itemize} \\
\end{tabularx}
\end{skillbox}

\begin{skillbox}[{Vocabulary, Concepts, \& Theorems}]{greenbox}
\renewcommand{\arraystretch}{1.4}
\begin{tabularx}{\linewidth}{@{} >{\leavevmode\bfseries\color{navy}}p{0.28\linewidth} X @{}}
Logarithmic model & A function of the form $f(x) = a\log_b x + k$ (or with a horizontal
                    shift) used to describe a situation where inputs change proportionally over
                    equal-length output intervals. \\
Proportion        & The constant factor by which the input value is multiplied as the output
                    increases by one unit in a logarithmic model. \\
Real zero         & The input value at which the logarithmic model output equals zero;
                    for $f(x) = \log_b x$, the real zero is $x = 1$. \\
Natural logarithm model & A logarithmic model using base $e$: $f(x) = a\ln(x) + k$, often
                    chosen for naturally-occurring phenomena. \\
Logarithmic regression & A technology-assisted technique that fits a logarithmic curve to a
                    data set by minimizing residuals. \\
Prediction        & Using a function model to estimate an output value for a given input
                    (or an input for a given output target). \\
\end{tabularx}
\end{skillbox}

\begin{skillbox}[Activate Prior Knowledge \& Spiral Review]{sky}
    Spiral review (warm-up) covers: reading an input-output table and deciding whether
    the pattern is linear, exponential, or logarithmic; identifying a real zero of a
    logarithmic function; and recognizing the ``proportional input'' pattern that signals
    a logarithmic model. These connect directly to EK 2.14.A.1 (recognizing the model type)
    and EK 2.14.A.2 (building a model from a zero and a proportion).
\end{skillbox}

\begin{skillbox}[Hook]{sky}
    Write on the board: \textbf{``A scientist notices that a population of bacteria changes at
    a decreasing rate: the population doubles in the first hour, doubles again after 2 more
    hours, then again after 4 more hours.''}

    Ask: ``What is happening to the \emph{time needed} for each doubling as the population
    grows?'' (It keeps increasing.) Then ask: ``If the output is the number of doublings and
    the input is the elapsed time, what kind of function might model this?'' Guide students
    to see that equal output increments (one doubling) correspond to multiplicatively growing
    input intervals (1 h, 2 h, 4 h --- each doubling). This motivates the logarithm as the
    right model when the \emph{input} changes proportionally for each unit of output.
\end{skillbox}

\begin{skillbox}[Lesson]{sky}
\begin{multicols}{2}
\textbf{Part 1 (10 min): When to Use a Logarithmic Model (EK 2.14.A.1).}

Present a table where the output increases by 1 each row, but the input multiplies by the
same factor each time. Ask: ``What does the output count?'' (How many times the input has
been multiplied by the proportion.) Contrast with an exponential model (equal input intervals,
proportional output). Have students label their notes table and identify the proportion $b$.

\columnbreak

\textbf{Part 2 (12 min): Building from Two Pairs (EK 2.14.A.2).}

Show the general form $f(x) = a\log_b x + k$. Use the worked example $f(1)=0$, $f(e)=3$:
since $\log_b 1 = 0$ for any $b$, we get $k = 0$; then $a\ln(e) = 3 \Rightarrow a = 3$.
Model the system-of-equations approach with a second example: $f(1)=2$, $f(9)=4$ with base 3.
Students fill in the guided-notes worked example.
\end{multicols}

\begin{multicols}{2}
\textbf{Part 3 (8 min): Transformations Approach (EK 2.14.A.3).}

Introduce $f(x) = a\log_b(x-h)+k$. Identify $a$ as vertical dilation, $h$ as horizontal
shift, $k$ as vertical shift. Present the decibel example as a real transformation. Students
match parameters to context descriptions.

\columnbreak

\textbf{Part 4 (10 min): Natural Log Models \& Predictions (EK 2.14.A.5, 2.14.A.6).}

Emphasize $\ln x = \log_e x$ and when it arises naturally. Show how to evaluate $f(e^n)$
and how to solve $a\ln x = c \Rightarrow x = e^{c/a}$. Students practice prediction in
context. Mention regression as a technology path (EK 2.14.A.4).
\end{multicols}
\end{skillbox}

\begin{skillbox}[Active Monitoring]{redbox}
\begin{itemize}[leftmargin=1.2em,itemsep=2pt,topsep=0pt]
  \item After Part 1: cold-call ``How does the pattern in this table differ from an exponential
        table?'' Expect: inputs multiply by a constant (not outputs).
  \item Watch for students who confuse the proportion (the base $b$) with the constant $a$.
        Emphasize: $b$ is identified from the proportional input pattern; $a$ is solved after.
  \item After Part 2: confirm students can set up two equations and solve the system correctly.
        Common error: forgetting that $\log_b 1 = 0$ always, giving $k$ for free.
  \item During Part 4: check that students can evaluate $f(e^2) = a\ln(e^2) = 2a$ without a
        calculator and can correctly invert $a\ln x = c$ to $x = e^{c/a}$.
\end{itemize}
\end{skillbox}

\begin{skillbox}[Group Work \& Differentiation]{redbox}
\begin{itemize}[leftmargin=1.2em,itemsep=2pt,topsep=0pt]
  \item \textbf{Tier R (Remediate):} Use a given table ($A = 1,10,100,1000$; $S = 0,1,2,3$)
        to identify the proportion, write $S(A) = \log_{10}(A)$, predict $S(10000)$, and
        discuss whether $S(0.5)$ is defined.
  \item \textbf{Tier A (Apply):} Given two input-output pairs for a species-diversity model,
        set up and solve a system for $a$ and $k$ in $f(x) = a\log_3 x + k$; write the
        complete model and make a prediction.
  \item \textbf{Tier E (Extension):} Analyze the decibel formula $S(I) = 10\log_{10}(I/I_0)$
        as a transformation; find the intensity of normal conversation (60 dB); compute the
        ratio of intensities between a rock concert (120 dB) and normal conversation.
\end{itemize}
\end{skillbox}

\begin{skillbox}[Individual Work \& Assessment]{redbox}
Exit ticket (2 questions, 1 page):
\begin{enumerate}[leftmargin=1.5em,itemsep=2pt,topsep=2pt]
  \item Given $x = 2, 4, 8, 16$ with $f(x) = 1, 2, 3, 4$: identify $b$, write
        $f(x) = \log_b x$, predict $f(32)$. (Tests EK 2.14.A.1 and 2.14.A.2.)
  \item Given $f(x) = 3\ln(x)$: evaluate $f(e^2)$ and $f(1)$; solve $f(x) = 6$ for $x$.
        (Tests EK 2.14.A.5 and 2.14.A.6.)
\end{enumerate}
Sort by result: students missing \#1 need work on recognizing proportional input patterns;
missing \#2 need practice evaluating and inverting natural log expressions.
\end{skillbox}

\begin{skillbox}[Reinforcement \& Extension]{goldbox}
\textbf{Homework (5--6 problems):} classify data sets as logarithmic vs.\ exponential;
construct $f(x) = a\ln x$ from two points; build $S(A) = a\log_{10} A + k$ from a four-row
table; analyze $f(x) = 2\log_2(x-1)+3$ (domain, evaluation, solving); AP-style MC on
context identification; extension on change-of-base connecting $a\ln x$ and $a\log_{10} x$.

\textbf{Preview of Lesson 2.15:} Connecting exponential and logarithmic models --- choosing
between them based on context, converting between forms, and using both in the same analysis.
\end{skillbox}

\end{document}
